\documentclass[11pt]{article}
\usepackage[T1]{fontenc}
\usepackage[margin=1in]{geometry}
\usepackage{amsmath,amssymb,amsthm}
\usepackage{enumitem}
\usepackage{graphicx}
\usepackage{microtype}
\usepackage{xcolor}
\usepackage[colorlinks=true,linkcolor=blue!55!black,urlcolor=blue!65!black]{hyperref}

\newtheorem{theorem}{Theorem}
\newtheorem{lemma}[theorem]{Lemma}
\newtheorem{corollary}[theorem]{Corollary}
\newcommand{\Commut}{\operatorname{Commut}}
\newcommand{\diam}{\operatorname{diam}}
\newcommand{\supp}{\operatorname{supp}}
\newcommand{\B}{\mathcal B}

\title{Property A and Operator-Norm Localization:\\
The Complete \(p\)-Phase Diagram}
\author{Candidate full answer to arXiv:1809.00532, Question 6.6}
\date{27 August 2026}

\begin{document}
\maketitle

\begin{center}
\fbox{\begin{minipage}{0.91\textwidth}
\textbf{Status: candidate full solution, likely valid.}
For every $1<p<\infty$, Proposition 6.3 remains true with $\ell^2(X)$
replaced by $\ell^p(X)$. At $p=1$, its localization clauses are universal and
do not characterize Property A.
\end{minipage}}
\end{center}

\section{Source question and definitions}

Špakula and Zhang prove a commutant version of operator norm localization that
characterizes Property A on $\ell^2(X)$, and ask on PDF page 20 of \cite{SZ}:

\begin{quote}
\textbf{Question 6.6.} Is it possible to generalise the characterisation of
Property A (Proposition 6.3) to any $p\in[1,\infty)$?
\end{quote}

Let $X$ be a bounded-geometry metric space and $1\le p<\infty$. For
$L,\eta>0$, set
\[
 \Commut_p(L,\eta)=\left\{b\in\B(\ell^p(X)):
 \|[b,f]\|<\eta\text{ for every $L$-Lipschitz }
 f\in C_b(X),\ \|f\|_\infty\le1\right\}.
\]
The direct $p$-versions of the three operator clauses in Proposition 6.3 are:

\begin{enumerate}[label=$(\mathrm{L}_{\arabic*}^p)$,leftmargin=3.1em]
\item For every $\varepsilon,L,M>0$ there is $s>0$ such that, if
$\|b\|\le M$ and $b\in\Commut_p(L,\varepsilon/6)$, some unit vector $v$ with
$\diam(\supp v)\le s$ satisfies $\|bv\|_p\ge\|b\|-\varepsilon$.

\item For every $c\in(0,1)$ and $L>0$ there is $s>0$ such that, if
$\|b\|=1$ and $b\in\Commut_p(L,(1-c)/6)$, some unit vector $v$ with
$\diam(\supp v)\le s$ satisfies $\|bv\|_p\ge c$.

\item ($p$-ONL) For every $c\in(0,1)$ and $R>0$ there is $s>0$ such that
every norm-one operator $a\in\B(\ell^p(X))$ of propagation at most $R$ has a
unit vector $v$ with $\diam(\supp v)\le s$ and $\|av\|_p\ge c$.
\end{enumerate}

For $E\subseteq X$, write
$\partial_R E=\{x\in X\setminus E:d(x,E)\le R\}$. Following \cite{ULA},
$X$ has \emph{uniform local amenability} (ULA) if for every $R,\varepsilon>0$
there is $S>0$ such that, for every finite $F\subseteq X$, some
$E\subseteq X$ with $\diam(E)\le S$ satisfies
\[
  |\partial_R E\cap F|<\varepsilon |E\cap F|.
  \tag{ULA}
\]

\section{Idea of the proof}

The source already proves Property A $\Rightarrow(\mathrm L_1^p)$ for every
$1<p<\infty$. The missing direction is geometric. Suppose ULA fails. At one
fixed scale $R$, arbitrarily large localization radii then admit a finite set
$F$ in which every small-diameter subset has a uniformly large outer vertex
boundary. Join points of $F$ at distance at most $R$ and consider the random
walk $P$ on this finite bounded-degree graph.

The lazy walk $(I+P)/2$ has norm one on every weighted $\ell^p$ because it
fixes constants. On the other hand, a coarea estimate turns local vertex
expansion into a local $\ell^2$ Cheeger inequality. The scalar monotonicity
inequality for $z\mapsto |z|^{p-2}z$ then gives a uniform lower bound for
$\|(I-P)v\|_p$ on every small-support vector. Uniform convexity of $\ell^p$
forces $(I+P)/2$ to lose a fixed amount of norm on all such vectors. Conjugating
from degree measure to counting measure produces a finite-propagation operator
that violates $p$-ONL. Thus $p$-ONL implies ULA, and Elek's theorem \cite{Elek}
then gives Property A.

At $p=1$, uniform convexity disappears for a decisive reason: every operator
is almost normed by a single column. Hence all three localization clauses are
automatic.

\section{The local lazy-walk estimate}

For a finite graph $G=(F,\mathcal E)$ with no isolated vertices, write $d_x$
for the degree of $x$ and
\[
 \|w\|_{p,d}^p=\sum_{x\in F}d_x|w_x|^p,
 \qquad
 (Pw)_x=\frac1{d_x}\sum_{y\sim x}w_y.
\]
Jensen's inequality shows that $P$ is a contraction on $\ell^p(F,d)$.

\begin{lemma}[Local Cheeger estimate]\label{lem:cheeger}
Suppose $G$ has maximum degree at most $D$, and $K\subseteq F$ has the
property that
\[
 |\partial_G A|\ge h|A|
 \tag{1}
\]
for every nonempty $A\subseteq K$, where $\partial_G A$ is the outer vertex
boundary in $F$. If $g\ge0$ is supported in $K$, then
\[
 \sum_{\{x,y\}\in\mathcal E}|g_x-g_y|^2
 \ge \frac{h^2}{2D^2}\sum_{x\in F}d_xg_x^2.
 \tag{2}
\]
\end{lemma}

\begin{proof}
For $A_t=\{x:g_x>t\}$, (1) implies that the number of edges leaving $A_t$ is
at least $h|A_t|$. The layer-cake formula therefore gives
\[
 \sum_{\{x,y\}\in\mathcal E}|g_x^2-g_y^2|
 =\int_0^\infty 2t\,|\mathcal E(A_t,A_t^c)|\,dt
 \ge h\sum_x g_x^2.
 \tag{3}
\]
By Cauchy--Schwarz and the degree bound, the left side of (3) is at most
\[
 \left(\sum_{\{x,y\}\in\mathcal E}|g_x-g_y|^2\right)^{1/2}
 \left(2D\sum_xg_x^2\right)^{1/2}.
\]
Squaring and using $\sum_xd_xg_x^2\le D\sum_xg_x^2$ proves (2).
\end{proof}

\begin{lemma}[Uniform local loss for the lazy walk]\label{lem:lazy}
Fix $1<p<\infty$, $D\ge1$, and $h>0$. There is $q=q(p,D,h)<1$ with the
following property. Under the hypotheses of Lemma~\ref{lem:cheeger}, every
$w\in\ell^p(F,d)$ supported in $K$ satisfies
\[
 \left\|\frac{I+P}{2}w\right\|_{p,d}\le q\|w\|_{p,d}.
 \tag{4}
\]
\end{lemma}

\begin{proof}
For $z,u\in\mathbb C$ and $1<p<\infty$, monotonicity of the duality map gives
\[
 \operatorname{Re}\!\left[(z-u)
 (|z|^{p-2}\overline z-|u|^{p-2}\overline u)\right]
 \ge \gamma_p\bigl(|z|^{p/2}-|u|^{p/2}\bigr)^2,
 \quad \gamma_p=\frac{4(p-1)}{p^2}.
 \tag{5}
\]
For complex scalars, the left side is minimized when $z$ and $u$ have the
same argument.  For $a\ge b\ge0$, Cauchy--Schwarz gives
\[
 \left(\int_b^a t^{p/2-1}\,dt\right)^2
 \le (a-b)\int_b^a t^{p-2}\,dt,
\]
which, after evaluating the two integrals, is exactly the resulting real
inequality in (5).

Pair $L=I-P$ against the $\ell^p(F,d)$ duality map. Symmetrizing over
unoriented edges, applying (5), and then Lemma~\ref{lem:cheeger} to
$g_x=|w_x|^{p/2}$ gives
\[
 \begin{split}
 \|Lw\|_{p,d}\|w\|_{p,d}^{p-1}
 &\ge \operatorname{Re}\sum_xd_x(Lw)_x|w_x|^{p-2}\overline{w_x}\\
 &\ge \gamma_p\sum_{\{x,y\}\in\mathcal E}
       (|w_x|^{p/2}-|w_y|^{p/2})^2\\
 &\ge \frac{\gamma_ph^2}{2D^2}\|w\|_{p,d}^p.
 \end{split}
 \tag{6}
\]
Thus $\|w-Pw\|_{p,d}\ge\delta\|w\|_{p,d}$, where
$\delta=\gamma_ph^2/(2D^2)>0$. Since $P$ is a contraction and
$\ell^p(F,d)$ has the same positive modulus of uniform convexity as an
$\ell^p$ space, (6) implies
\[
 \left\|\frac{w+Pw}{2}\right\|_{p,d}
 \le (1-\rho_p(\delta))\|w\|_{p,d},
\]
where $\rho_p(\delta)>0$ depends only on $p$ and $\delta$. This is (4).
\end{proof}

\section{The new reverse implication}

\begin{theorem}\label{thm:onl-ula}
Let $X$ have bounded geometry and let $1<p<\infty$. If $X$ has $p$-ONL,
then $X$ has ULA.
\end{theorem}

\begin{proof}
Suppose ULA fails. Then there are $R>0$ and $h>0$ such that for every $S>0$
there is a finite $F\subseteq X$ satisfying
\[
 |\partial_R E\cap F|\ge h|E\cap F|
 \tag{7}
\]
for every $E\subseteq X$ with $\diam(E)\le S$. Decreasing $h$ if necessary,
assume $h\le1$.

On $F$, form the graph $G_F$ by joining distinct $x,y$ when $d(x,y)\le R$.
Its maximum degree is at most
$D=\sup_x|B(x,R)|-1$. Taking $E=\{x\}$ in (7) shows that $G_F$ has no
isolated vertices. If $A\subseteq F$ has diameter at most $S$, then
$\partial_{G_F}A=\partial_RA\cap F$, so (7) gives (1).

Let $P_F$ be the random walk of $G_F$ and $A_F=(I+P_F)/2$ on
$\ell^p(F,d)$. Lemma~\ref{lem:lazy} supplies a number $q<1$, depending only
on $p,R,h$ and the bounded-geometry bound, such that
\[
 \|A_Fw\|_{p,d}\le q\|w\|_{p,d}
 \tag{8}
\]
whenever $\diam(\supp w)\le S$.

The isometry $U_F:\ell^p(F)\to\ell^p(F,d)$ defined by
$(U_Fv)_x=d_x^{-1/p}v_x$ preserves supports. Define on $\ell^p(X)$
\[
 T_F=U_F^{-1}A_FU_F\quad\text{on }\ell^p(F),
 \qquad T_F=0\quad\text{on }\ell^p(X\setminus F).
\]
The operator $T_F$ has propagation at most $R$ and norm one: $A_F$ is a
contraction and fixes the constant function. For every $v$ of support diameter
at most $S$, (8) applied to $U_F(v|_F)$ yields
\[
             \|T_Fv\|_p\le q\|v|_F\|_p\le q\|v\|_p.
 \tag{9}
\]
Choose $c\in(q,1)$. For the fixed propagation $R$, equation (9) gives, for
every proposed localization radius $S$, a norm-one operator $T_F$ for which
no $S$-localized unit vector reaches $c$. This contradicts $p$-ONL.
\end{proof}

\begin{corollary}\label{cor:onl-a}
For a bounded-geometry metric space and $1<p<\infty$, $p$-ONL implies
Property A.
\end{corollary}

\begin{proof}
By Theorem~\ref{thm:onl-ula}, $X$ has ULA. We record why Elek's graph theorem
\cite{Elek} applies without a hidden coarse-geodesicity assumption. Fix $R>0$
and let $G_R$ have vertex set $X$ and edges between distinct points at distance
at most $R$. Its degree is bounded by some $D$. Given $\varepsilon>0$ and a
finite subgraph $H$ on a vertex set $F$, metric ULA at scale $R$, with tolerance
$\varepsilon/D$, produces $E$ of uniformly bounded diameter with
$|\partial_RE\cap F|<(\varepsilon/D)|E\cap F|$. Put $A=E\cap F$. Bounded
geometry uniformly bounds $|A|$. Every inner boundary vertex of $A$ in $H$
has a neighbor in $\partial_RE\cap F$, and each such outer vertex is charged
at most $D$ times. Hence $|\partial_H^{\rm in}A|<\varepsilon|A|$, so the class
of finite subgraphs of $G_R$ is uniformly locally amenable. Elek's theorem
gives Property A for $G_R$. Its controlled supports in graph distance are
controlled in $X$-distance, and every pair at $X$-distance at most $R$ is an
edge of $G_R$. Applying this for each $R$ gives Property A for $X$ (the finite
case being immediate).
\end{proof}

\section{Completion of the characterization}

\begin{lemma}[Finite propagation implies slow commutation]\label{lem:comm}
For every $R,\eta>0$ there is $L>0$ such that every norm-one operator on
$\ell^p(X)$ of propagation at most $R$ belongs to $\Commut_p(L,\eta)$.
\end{lemma}

\begin{proof}
Bounded geometry lets one decompose the controlled relation
$\{(x,y):d(x,y)\le R\}$ into $N=N(X,R)$ graphs of partial translations.
Accordingly, every such operator has a decomposition
$a=\sum_{k=1}^N M_{g_k}V_k$, where $V_k$ is a partial translation of
propagation at most $R$ and $\|g_k\|_\infty\le1$; this is the same support
decomposition used in Lemma 6.2 of \cite{SZ}, and it is independent of $p$.
If $f$ is an $L$-Lipschitz contraction, then
\[
 \|[M_{g_k}V_k,f]\|\le LR,
 \qquad
 \|[a,f]\|\le NLR.
\]
Take $L<\eta/(NR)$ (with the propagation-zero case immediate).
\end{proof}

\begin{theorem}[Complete answer to Question 6.6]\label{thm:main}
Let $X$ be a bounded-geometry metric space.

\begin{enumerate}[label=\textup{(\roman*)}]
\item For every fixed $1<p<\infty$, the following are equivalent:
Property A; the metric sparsification property; and
$(\mathrm L_1^p)$, $(\mathrm L_2^p)$, $(\mathrm L_3^p)$ above.

\item At $p=1$, the three localization conditions
$(\mathrm L_1^1)$, $(\mathrm L_2^1)$, $(\mathrm L_3^1)$ hold for every metric
space, so none of them characterizes Property A.
\end{enumerate}
\end{theorem}

\begin{proof}
Fix $1<p<\infty$. Property A is equivalent to metric sparsification, and
metric sparsification implies $(\mathrm L_1^p)$ by the $p$-version of the
commutant localization lemma proved in Section 5 of \cite{SZ}. The implication
$(\mathrm L_1^p)\Rightarrow(\mathrm L_2^p)$ follows by taking
$M=1$ and $\varepsilon=1-c$. Lemma~\ref{lem:comm} gives
$(\mathrm L_2^p)\Rightarrow(\mathrm L_3^p)$. Finally,
Corollary~\ref{cor:onl-a} gives $(\mathrm L_3^p)\Rightarrow$ Property A.

For $p=1$, every $b\in\B(\ell^1(X))$ satisfies the exact column formula
\[
                 \|b\|=\sup_{x\in X}\|b\delta_x\|_1.
 \tag{10}
\]
Indeed, the reverse inequality follows from
$\|bv\|_1\le\sum_x|v_x|\|b\delta_x\|_1$; the other inequality is immediate.
Thus a one-point-supported unit vector $\varepsilon$-norms every $b$, and it
exceeds any fixed $c<1$ when $\|b\|=1$. All three localization clauses hold
with localization diameter zero, without their commutator or propagation
hypotheses. A coarse disjoint union of bounded-degree expanders has bounded
geometry but not Property A, proving failure of characterization at $p=1$.
\end{proof}

\section{Novelty, limitations, and verification}

Theorem~\ref{thm:main} answers the direct formulation of Question 6.6 across
its full stated range. It does not claim that no differently designed endpoint
condition could characterize Property A at $p=1$.

Chung and Nowak \cite{CN} prove that Property A implies $p$-ONL for every
$p\in[1,\infty)$ and explicitly state that they do not know whether the
converse holds, suggesting ULA as a possible route. A bounded search used the
exact source question, ``$p$-operator norm localization'', ``independent of
$p$'', ``Property A converse'', ``uniform local amenability'', the cited-paper
graph, and the run's cheap indexes. No later resolution through 27 August 2026
was located. Priority is therefore plausible, not certified.

The proof has no computational dependency. The main checks are: the outer
boundary in failed ULA is exactly the graph boundary used in (7); all constants
in Lemmas~\ref{lem:cheeger}--\ref{lem:lazy} depend only on $p,R,h$ and bounded
geometry, never on $F$ or $S$; degree conjugation is an isometry and preserves
supports; and the source's forward commutant lemma already covers every
$1<p<\infty$. Human review should concentrate on these four points and the
metric-to-graph application of Elek's theorem.

\clearpage
\section*{Source-page evidence}
\begin{center}
  \includegraphics[width=\textwidth,height=0.9\textheight,keepaspectratio]
    {figures/open_problem_crop.png}
\end{center}

\begin{thebibliography}{9}
\bibitem{SZ}
J.~Špakula and J.~Zhang,
\emph{Quasi-Locality and Property A},
J. Funct. Anal. 278 (2020), 108299; arXiv:1809.00532.

\bibitem{ULA}
J.~Brodzki, G.~A. Niblo, J.~Špakula, R.~Willett, and N.~Wright,
\emph{Uniform local amenability},
J. Noncommut. Geom. 7 (2013), 583--603; arXiv:1203.6169.

\bibitem{Elek}
G.~Elek,
\emph{Uniform Local Amenability implies Property A},
Proc. Amer. Math. Soc. 149 (2021), 2573--2577; arXiv:1912.00806.

\bibitem{CN}
Y.~C. Chung and P.~W. Nowak,
\emph{Expanders are counterexamples to the $\ell^p$ coarse Baum--Connes
conjecture},
J. Noncommut. Geom. 17 (2023), 305--331; arXiv:1811.10457.
\end{thebibliography}

\end{document}
